\section{Uvod}

Mali modularni reaktori (SMR) predstavljaju jednu od tehnologija koja se u raspravama o budućnosti nuklearne energetike često povezuje s modularnošću, kraćim vremenom izgradnje, sigurnosnim poboljšanjima i potencijalno fleksibilnijim uklapanjem u elektroenergetski sustav. Uz tehničke i ekonomske argumente, za prihvat novih energetskih tehnologija važna je i percepcija javnosti.

Cilj ovog izvještaja je analizirati odgovore ispitanika iz ankete provedene utorkom, s naglaskom na pitanja dodijeljena Timu E. Analiza se usredotočuje na percepciju najveće prednosti SMR tehnologije te na stav prema državnom subvencioniranju razvoja SMR tehnologije umjesto ulaganja u jednu klasičnu veliku nuklearnu elektranu.

Iako naslov rada obuhvaća širi kontekst sigurnosti i prihvatljivosti izgradnje, statistička obrada u ovom izvještaju temelji se na pitanjima P9 i P10, uz korištenje općih demografskih podataka o ispitanicima.
